\chapter{Theoretical Framework and Methodology}
\label{chap:theory}
\label{chap:methodology}

\section{Empirical Objective}

The thesis estimates horizon-specific aggregate responses to unexpected ECB policy news. The canonical shock is the monthly high-frequency target surprise, signed so positive values denote easing. Responses cover housing credit, lending conditions, financial-market prices, liquidity-stock context, broad credit quantities, labor expectations, and compensation-linked wage-pressure measures. The estimand is comparative: which channels show the clearest and most persistent response profiles?

The core evidence is therefore a path, not a single impact coefficient. Signs, accumulated responses, uncertainty bands, and stability diagnostics are read together. The model is presented as reduced-form dynamic evidence on heterogeneous monetary transmission, not as a structural welfare analysis or household distribution study.

\section{Mechanism Framework}

The interpretation is organized around a housing-credit transmission chain:
\begin{equation}
\begin{aligned}
  \text{ECB policy news}
  &\rightarrow
  \text{intermediation variables} \\
  &\rightarrow
  \text{housing-credit expansion}
  \rightarrow
  \text{selected financial responses}.
\end{aligned}
\label{eq:mechanism-chain}
\end{equation}

The first link is measured by high-frequency ECB announcement surprises. The second link is represented by lending spreads, lending rates, broad credit aggregates, liquidity-stock context, and bank-related timing evidence. The third link is represented primarily by house-purchase lending growth and secondarily by pure-new mortgage flows. Equity and broader financial variables support the macro-financial interpretation but do not define the thesis on their own.

This chain is not a structural mediation model. The variables respond to the same policy-news shock, and aggregate data cannot isolate bank-level decisions or household incidence. The framework instead provides an economically disciplined way to read timing and persistence: intermediation variables should move in ways that make the housing-credit response plausible, while compensation-linked variables provide the comparative real-economy block.

\section{Identification Problem in Monetary Economics}

Monetary policy is endogenous to inflation, slack, financial stress, and central-bank forecasts. At monthly or quarterly frequency, a policy-rate movement can mix the action itself, private-sector anticipation, and information released by the central bank. A generic reduced-form VAR illustrates the limitation:
\begin{equation}
  \vect{y}_t
  =
  \vect{c}
  +
  \sum_{j=1}^{p}\mat{A}_j\vect{y}_{t-j}
  +
  \vect{u}_t,
  \qquad
  \Expect[\vect{u}_t\vect{u}_t\transpose]=\mat{\Sigma}_u .
  \label{eq:reduced-form-var}
\end{equation}
Without additional restrictions, \(\vect{u}_t\) does not separate unexpected policy news from macroeconomic information.

High-frequency identification improves timing by measuring yield or asset-price changes inside narrow ECB announcement windows \parencite{Kuttner2001,GertlerKaradi2015,NakamuraSteinsson2018,Altavilla2019}. It does not remove information effects, because announcements can reveal the central bank's assessment of future conditions \parencite{JarocinskiKaradi2020}. The surprise is therefore used as measured policy news, not as a fully purified structural innovation.

\section{Construction of Monetary Policy Shock Series}

The shock pipeline is implemented in \nolinkurl{src/data/ecb_monetary_surprises.py} and consumed by \nolinkurl{src/data/build_monthly_model_dataset.R}. For event \(e\) in month \(m\), let \(z^{f}_{e}\) denote the high-frequency surprise for factor family \(f\). The monthly factor is
\begin{equation}
  Z^{f}_{m}
  =
  \sum_{e \in \mathcal{E}(m)} z^{f}_{e},
  \label{eq:monthly-shock-aggregation}
\end{equation}
where \(\mathcal{E}(m)\) is the set of ECB events in month \(m\). No-event months receive zero values inside source coverage; post-source months are excluded.

The thesis-facing shock is
\begin{equation}
  \mathrm{Shock}_{m}
  =
  - Z^{target}_{m},
  \label{eq:target-easing-sign}
\end{equation}
stored as \nolinkurl{target_factor_monthly_easing}. The sign flip makes positive responses correspond to easing surprises. Timing, forward-guidance, QE, and weighted-composite variants enter robustness checks.

\section{Information-Effect Screening}

The information-effect screen compares the screening rate factor, the OIS six-month response, and the STOXX50 response. Opposite yield-equity signs are classified as cleaner monetary events; same-sign, conflicting, zero, or missing windows are treated as contaminated or ambiguous. In the current monthly identification window, 76 events are clean and 132 are contaminated or ambiguous. The baseline keeps the full target shock, while clean-event and contaminated-event series test sensitivity.

The event-density diagnostics in \Cref{fig:instrument-quarter-density,fig:instrument-regime-histograms} document the calendar concentration and regime composition behind these screens. The DAX response is interpreted under this same identification caution: an easing surprise may include adverse macroeconomic information communicated by the ECB.

\section{Relevance Diagnostics and Interpretation Boundary}

The repository tests whether high-frequency surprises strongly bridge to low-frequency treatment variables such as DFR changes, the Wu-Xia shadow rate, or ECB assets \parencite{MertensRavn2013,StockWatson2018}. The tournament in \nolinkurl{scripts/run_first_stage_tournament.py} records \(F\)-statistics, partial \(R^2\), event concentration, jackknife sensitivity, rolling relevance, sign stability, and regime strength \parencite{StockYogo2005}.

The strongest bridge is the end-of-month DFR change on \nolinkurl{target_factor_monthly_easing}, with \(F=17.07\) and partial \(R^2=0.066\). This is meaningful full-sample relevance rather than a weak-instrument collapse. The caution is about sensitivity: the jackknife \(F\)-ratio falls to about 0.31 after removing the largest shocks, and the rolling median \(F\)-statistic is about 8.99. The final design therefore estimates responses to policy news itself instead of promoting a stronger low-frequency treatment interpretation.

\section{Local Projection Framework}

The active estimator is a monthly local projection \parencite{Jorda2005}. For response \(y\), horizon \(h\), and shock \(\mathrm{Shock}_{t}\),
\begin{equation}
  y_{t+h} - y_{t-1}
  =
  \alpha_h
  + \beta_h \mathrm{Shock}_{t}
  + \Gamma_h X_t
  + \varepsilon_{t+h},
  \qquad
  h \in \{0,1,3,6,12,24\}.
  \label{eq:baseline-lp}
\end{equation}
\(X_t\) includes two lags of the response, two lags of euro-area monthly inflation, and two lags of the monthly-average DFR where available.

Responses are normalized to a one-standard-deviation shock:
\begin{equation}
  \widehat{\IRF}_{h}
  =
  \widehat{\beta}_{h}\widehat{\sigma}_{\mathrm{Shock}}.
  \label{eq:lp-normalization}
\end{equation}
Accumulated paths provide the main duration comparison.

\section{Uncertainty Quantification}

Baseline inference uses HAC/Newey--West covariance estimates with bandwidth \(h+1\), implemented in \nolinkurl{src/lpiv/ols_hac.py} \parencite{NeweyWest1987}. The uncertainty layer also reports circular block, moving-block, and wild-bootstrap variants under \nolinkurl{results/final/uncertainty/}. These checks govern long-horizon claims where adjacent LP estimates share innovations.

\section{Sequential Timing Analysis}

The sequential layer asks whether earlier movements in lending conditions, liquidity, credit, or labor expectations are followed by later movements in downstream variables. For \(h_u<h_d\),
\begin{equation}
  d_{t+h_d} - d_{t-1}
  =
  a
  + \lambda (u_{t+h_u} - u_{t-1})
  + \theta \mathrm{Shock}_{t}
  + \Pi W_t
  + \nu_t .
  \label{eq:sequential-timing}
\end{equation}
\(\lambda\) is an ordered timing association. The evidence remains timing-consistent rather than structurally mediated.

\section{What the Model Is and Is Not}

The model is reduced-form dynamic causal evidence under imperfect identification. It is a heterogeneous transmission analysis and a macro-financial propagation study centered on post-COVID-relevant ECB policy news. It is not a structural welfare analysis, an inequality decomposition, a household distribution study, or evidence on redistribution. This boundary governs the results chapter and the conclusion.

\section{Robustness Governance}

The pipeline \nolinkurl{scripts/run_full_pipeline.py} rebuilds data, shocks, diagnostics, proxy validation, LP outputs, uncertainty checks, timing tables, figures, memos, and the final audit. The active results surface is \nolinkurl{results/final/}. The stress tests cover event contamination, historical exclusions, rolling and recursive windows, alternative shock factors, and descriptive regimes. The next chapter documents the data and proxy architecture.
